% Build with: latexmk -xelatex JEVANY_METHOD_AND_ABLATIONS.tex
\documentclass[10pt]{article}

\usepackage{fontspec}
\setmainfont{DejaVu Sans}
\setsansfont{DejaVu Sans}
\setmonofont{DejaVu Sans Mono}
\usepackage{unicode-math}
\setmathfont{Latin Modern Math}
\usepackage[letterpaper,margin=0.72in,headheight=15pt]{geometry}
\usepackage{booktabs,array}
\usepackage{xcolor}
\usepackage{hyperref}

\definecolor{navy}{HTML}{101A35}
\definecolor{blue}{HTML}{2563EB}
\definecolor{cyan}{HTML}{06B6D4}
\definecolor{violet}{HTML}{7C3AED}
\definecolor{mint}{HTML}{10B981}
\definecolor{amber}{HTML}{F59E0B}
\definecolor{ink}{HTML}{172033}
\definecolor{muted}{HTML}{5F6B7A}
\definecolor{paper}{HTML}{F7F9FC}
\definecolor{line}{HTML}{D9E1EC}
\definecolor{softblue}{HTML}{EEF4FF}
\definecolor{softmint}{HTML}{ECFDF5}
\definecolor{softamber}{HTML}{FFF7E6}

\hypersetup{
  colorlinks=true,
  linkcolor=blue,
  urlcolor=violet,
  pdftitle={JevAny: Method, Training, and Ablation Report},
  pdfauthor={JevAny contributors}
}

\setlength{\parindent}{0pt}
\setlength{\parskip}{5pt}
\renewcommand{\arraystretch}{1.16}

\makeatletter
\renewcommand\section{\@startsection{section}{1}{\z@}%
  {-3.0ex \@plus -1ex \@minus -.2ex}{1.4ex \@plus .2ex}%
  {\color{navy}\Large\bfseries}}
\renewcommand\subsection{\@startsection{subsection}{2}{\z@}%
  {-2.4ex \@plus -1ex \@minus -.2ex}{.8ex \@plus .2ex}%
  {\color{blue}\large\bfseries}}
\makeatother

\pagestyle{plain}

\newcommand{\good}[1]{\textcolor{mint!60!black}{\textbf{#1}}}
\newenvironment{summarybox}{%
  \par\medskip\noindent\begin{center}\begin{minipage}{.94\linewidth}%
  \color{ink}\textcolor{blue}{\rule{3pt}{\baselineskip}}\hspace{7pt}%
}{\end{minipage}\end{center}\medskip}

\begin{document}

\begin{titlepage}
\pagecolor{navy}\color{white}

\vspace*{1.30in}
{\color{cyan}\rule{\textwidth}{5pt}\par}\vspace{0.32in}
{\color{white}\fontsize{36}{39}\selectfont\bfseries JevAny\par}
\vspace{0.16in}
{\color{white}\fontsize{21}{25}\selectfont Method, Training, and Ablation Report\par}
\vspace{0.22in}
{\color{cyan!35}\large One decision interface across open language-model backbones\par}

\vfill
\noindent\fcolorbox{white!22}{navy}{\begin{minipage}{.78\textwidth}\color{white}\vspace{5pt}
\textbf{Release snapshot} \hfill 30 September 2026\\[4pt]
Five open LoRA checkpoints $\cdot$ Pointer and direct-token readouts $\cdot$
Transfer and JevBench evaluation $\cdot$ Reproducible training telemetry
\vspace{5pt}\end{minipage}}\par
\vspace{0.22in}
{\color{white!68}\small
Repository: \href{https://github.com/SimpleJev/JevAny}{github.com/SimpleJev/JevAny}\\
Models: \href{https://huggingface.co/collections/tianxinwei/jevany-adaptive-decision-systems-6ab2c941bcecb4d2c61d1326}{JevAny model collection}\par}
\end{titlepage}
\nopagecolor\color{ink}

\tableofcontents
\clearpage

\section{Executive summary}

\begin{summarybox}
\textbf{JevAny turns a language-model backbone into a bounded decision model.}
It consumes a shared state, one or more questions, and explicit candidate answers;
it returns a normalized distribution over the supplied candidates without decoding
answer text. The current release trains rank-8 LoRA adapters in BF16 with a
cross-entropy objective and either a residual pointer head or a direct-token
readout. The strongest released model, Qwen3.8 27B LoRA, reaches
\textbf{85.76\% Transfer accuracy} and \textbf{90.48\% JevBench accuracy}.
\end{summarybox}

The release family spans Gemma 4B, Qwen3.5 4B, Qwen3.8 27B, and Muse Glimmer
30B. The training corpus contains \textbf{1,772,725 text records} and
\textbf{2,180,242 labelled decisions} across preference, agent/tool decision,
reasoning, classification, and safety tasks. This report describes the corpus
scale and task families.

The ablation program used
fixed-step screens, held-out calibration, paired record-group bootstrap intervals,
explicit promotion gates, and full-run confirmation. Pure InfoNCE, supervised
Brier regularization, RLCR variants, targeted mixture changes, and dual-tower
alternatives failed their predeclared promotion gates.

\section{Problem formulation}

For a state $s$, question $q$, and explicit choices
$\mathcal{C}=\{c_1,\ldots,c_K\}$, JevAny predicts
\[
  p(y=i\mid s,q,\mathcal{C}), \qquad i\in\{1,\ldots,K\}.
\]
The candidate set is part of the model input, and the output is already a
probability simplex suitable for ranking, confidence thresholds, or downstream
control.

\begin{center}
\fcolorbox{line}{paper}{\parbox[c][12mm][c]{25mm}{\centering\bfseries State +\\question + choices}}
$\;\longrightarrow\;$
\fcolorbox{blue!25}{softblue}{\parbox[c][12mm][c]{25mm}{\centering\bfseries Structured\\causal encoding}}
$\;\longrightarrow\;$
\fcolorbox{blue!25}{softblue}{\parbox[c][12mm][c]{22mm}{\centering\bfseries Frozen LM\\+ LoRA}}
$\;\longrightarrow\;$
\fcolorbox{mint!30}{softmint}{\parbox[c][12mm][c]{25mm}{\centering\bfseries Decision\\readout}}
$\;\longrightarrow\;$
\fcolorbox{mint!30}{softmint}{\parbox[c][12mm][c]{20mm}{\centering\bfseries $p(c_i)$}}
\end{center}

\subsection{Structured causal input}

The text layout is
\[
\langle\mathrm{state}\rangle\;s\;
\langle\mathrm{question}\rangle\;q\;
\prod_{i=1}^{K}\bigl(\langle\mathrm{option}\rangle\;c_i\;
\langle/\mathrm{option}\rangle\bigr)\;
\langle\mathrm{decide}\rangle.
\]
User text is escaped so it cannot forge control delimiters. When a backbone has
validated custom-mask support, questions may be packed with a block-causal mask.
Recurrent, sliding-window, and other unsupported architectures instead receive
independent causal rows that repeat the shared state. This preserves
question-level isolation across backbone families.

\subsection{Pointer readout}

Let $h_d$ be the hidden state at the decision marker and $h_i$ the state at the
closing marker of option $i$. With pointer width $d_p=256$,
\[
 q_d=W_qh_d+b_q,\qquad k_i=W_kh_i+b_k,\qquad
 z_i=\frac{k_i^\top q_d}{\sqrt{d_p}}.
\]
The released pointer checkpoints add a zero-output-initialized residual scorer,
\[
 z_i\leftarrow z_i+w_2^\top\operatorname{GELU}(W_1(k_i\odot q_d)+b_1)+b_2,
\]
where the residual width is 256. Softmax is applied only across the options in
the request. Consequently, the head has no fixed option-vocabulary ceiling; the
practical choice limit is the model context window.

\subsection{Direct-token readout}

Direct-token mode prefixes candidates with a deterministic table of distinct,
round-tripping single-token labels. At the decision marker, the frozen base LM
head produces full-vocabulary logits $v$ and training minimizes
\[
  \mathcal{L}_{\mathrm{DT}}=-\sum_{i=1}^{K}t_i
  \log\operatorname{softmax}(v)_{\nu_i},
\]
where $\nu_i$ is the vocabulary ID assigned to option $i$. Inference selects
the $K$ candidate logits and renormalizes them. The released implementation
supports up to 255 candidates.

\begin{summarybox}
\textbf{Pointer vs direct-token at a glance}\par\smallskip
\begin{tabular}{>{\bfseries}p{.18\linewidth}p{.36\linewidth}p{.36\linewidth}}
\toprule
 & Pointer & Direct-token \\
\midrule
Scoring & Learned query--option compatibility & Frozen LM head over fixed label tokens \\
Training & Candidate-only softmax; efficient & Full-vocabulary normalization; slower \\
Choices & More than 255, context permitting & Up to 255 \\
Inference & One backbone prefill, no decoding & One backbone prefill, no decoding \\
Observed 4B result & Better Transfer: 78.68\% & Better JevBench: 80.95\% \\
\bottomrule
\end{tabular}
\end{summarybox}

Both paths have similar expected inference latency for comparable inputs because
neither autoregressively generates an answer. In the complete Qwen 4B training
trajectory, direct-token consumed 381.9 H200 GPU-hours versus 330.6 for pointer,
about 15.5\% more; the released direct-token checkpoint was selected earlier in
that trajectory.

\section{Training method}

\subsection{Supervised objective}

For a hard label $y$, the default objective is
\[
  \mathcal{L}_{\mathrm{CE}}=-\log p_y.
\]
For a normalized soft target $t$, it becomes
\[
  \mathcal{L}_{\mathrm{soft}}=-\sum_i t_i\log p_i.
\]
The release runs use CE, rank-8 LoRA over the backbone's supported linear
projections, zero LoRA dropout, a one-cycle learning-rate schedule, BF16 weights
and autocast, and seed 0. Pointer models train the 256-dimensional residual head
alongside LoRA. Direct-token freezes the original LM head and trains only the
adapter path.

\subsection{Optional research objectives}

The training code also supports contrastive terms, supervised Brier penalties,
and an experimental calibration-aware policy stage. These were ablated, but are
not part of the five current releases.

For RLCR, noisy logit proposals $z'=z+\epsilon$ use
$\epsilon\sim\mathcal{N}(0,\sigma^2I)$ after per-proposal mean centering. If
$a=\arg\max_i\operatorname{softmax}(z')_i$, confidence is
$q=\operatorname{softmax}(z')_a$, correctness is $c=t_a$, and reward is
\[
  r(c,q)=c-(q-c)^2.
\]
A group-relative advantage multiplies the Gaussian location score, while a CE
anchor limits decision drift. Policy optimization operates directly on perturbed
decision logits, without generating reasoning rollouts.

\subsection{Production release recipe}

\begin{tabular}{p{.30\linewidth}rrrrr}
\toprule
Model & Readout & LoRA & Global batch & GPUs & Released step \\
\midrule
Gemma 4B & Pointer & 8 & 384 & 32 & 2,771 \\
Qwen3.5 4B & Pointer & 8 & 128 & 32 & 13,850 \\
Qwen3.5 4B & Direct-token & 8 & 128 & 32 & 9,695 \\
Muse Glimmer 30B & Pointer & 8 & 160 & 40 & 3,324 \\
Qwen3.8 27B & Pointer & 8 & 32 & 32 & 22,160 \\
\bottomrule
\end{tabular}

All five runs use one epoch, learning rate $5\times10^{-5}$, weight decay 0.01,
BF16 training, gradient checkpointing, and LoRA targets resolved from the loaded
architecture. The release data scale is fixed at 1,772,725 text records and
2,180,242 labelled decisions.

\section{Evaluation protocol}

\subsection{Suites}

\begin{itemize}
  \item \textbf{Transfer} contains 1,046 clean, knowable decisions spanning seven
  cross-domain datasets and four robustness slices. Every reported run covers all
  1,046 scored items.
  \item \textbf{JevBench} contains all 231 items in the frozen public development
  conversion, yielding public-development scores.
\end{itemize}

Transfer accuracy is an item-level mean. NLL, multiclass Brier score, and
10-bin ECE use the candidate probabilities. A scalar temperature is fit on a
separate calibration partition by micro NLL and then applied unchanged to the
evaluation rows. Temperature scaling cannot change argmax accuracy.

For controlled ablations, comparisons use identical rows and option order.
Uncertainty is estimated by paired, source-stratified record-group bootstrap;
sibling questions remain together. Registered studies use 2,000 or 5,000
resamples and apply explicit non-inferiority and calibration-improvement gates.

\subsection{Released model results}

\begin{tabular}{p{.32\linewidth}rrrrr}
\toprule
Model & Transfer $\uparrow$ & JevBench $\uparrow$ & NLL $\downarrow$ & Brier $\downarrow$ & ECE $\downarrow$ \\
\midrule
Gemma 4B LoRA & 70.84 & 77.49 & 0.706 & 0.369 & 0.056 \\
Qwen3.5 4B LoRA & 78.68 & 80.09 & 0.587 & 0.297 & 0.035 \\
Qwen3.5 4B Direct-Token & 78.20 & 80.95 & 0.564 & 0.291 & \good{0.029} \\
Muse Glimmer 30B LoRA & 83.46 & 87.45 & 0.464 & 0.229 & 0.032 \\
Qwen3.8 27B LoRA & \good{85.76} & \good{90.48} & \good{0.392} & \good{0.200} & 0.030 \\
\bottomrule
\end{tabular}

All accuracy values are percentages. NLL, Brier, and ECE are Transfer metrics.
Unrounded values and suite hashes live in
\href{https://github.com/SimpleJev/JevAny/blob/main/results/model-family-v2.json}{the machine-readable release record}.

\section{Ablation program}

\subsection{Ablation procedure}

The ablation program screened candidates at fixed budgets and confirmed them
at full scale:

\begin{enumerate}
  \item \textbf{Register} the arm, frozen inputs, causal comparison step, metrics,
  and promotion gate before training.
  \item \textbf{Audit step zero} for identical predictions or initialized weights
  whenever an arm reuses a parent.
  \item \textbf{Screen at fixed budget} so capacity, head, readout, and objective
  candidates receive matched exposure.
  \item \textbf{Recalibrate independently} on held-out calibration rows; never fit
  temperature on Transfer or JevBench.
  \item \textbf{Compare paired rows} with group bootstrap intervals and reject an
  arm if either accuracy non-inferiority or calibration improvement fails.
  \item \textbf{Confirm at scale} before promoting a screen winner to the
  release recipe.
\end{enumerate}

The early 4B sweep favored rank 32, query-mean, and a
small contrastive weight by aggregate rank at its fixed causal checkpoint. The
separate full model-family runs use rank 8, decision-marker/option-close, and CE,
which is the configuration recorded by the published checkpoints.

\subsection{Run validity and metric corrections}

One CE-path isolation run was
marked invalid after its parity audit and is excluded from all conclusions.
Capacity probes, failed launches, and telemetry-recovery jobs likewise do not
enter the result tables. A post-terminal audit also found that two early reviews
had reported raw rather than temperature-scaled NLL because saved logits took
precedence over calibrated probabilities in the metric helper. The amended NLL
values are used here; accuracy, Brier, ECE, promotion gates, and selected arms
were verified unchanged.

\subsection{Pointer-head structure}

A controlled 250-step screen compared the dot-product head, an MLP scorer, and
a zero-initialized residual correction. The table shows calibrated development
metrics over 904 decisions.

\begin{tabular}{p{.29\linewidth}rrrrr}
\toprule
Head & Accuracy $\uparrow$ & NLL $\downarrow$ & Brier $\downarrow$ & ECE $\downarrow$ & Parameters \\
\midrule
Linear pointer & 81.97 & 0.4776 & 0.2603 & 0.0671 & 2,621,952 \\
MLP scorer & 82.41 & 0.4935 & 0.2596 & 0.0722 & 2,688,001 \\
Residual pointer & \good{82.63} & \good{0.4586} & \good{0.2497} & \good{0.0601} & 2,688,001 \\
\bottomrule
\end{tabular}

The residual arm passed the registered gate and was selected. Its final layer
starts at zero, preserving the base pointer score at initialization while
allowing a nonlinear correction during training.

\subsection{Representation readout}

The early Qwen 4B study compared the decision marker and option closing marker
against mean-pooled query, option, and combined representations. Results below
are the fixed-screen development values.

\begin{tabular}{p{.38\linewidth}rrrr}
\toprule
Readout & Accuracy $\uparrow$ & NLL $\downarrow$ & Brier $\downarrow$ & ECE $\downarrow$ \\
\midrule
Decision / option-close control & 77.88 & 0.5405 & 0.3065 & 0.0642 \\
Query-mean & \good{78.87} & \good{0.5117} & \good{0.2991} & \good{0.0428} \\
Option-mean & 74.45 & 0.6422 & 0.3653 & 0.0514 \\
Combined means & 73.34 & 0.6322 & 0.3709 & 0.0628 \\
\bottomrule
\end{tabular}

Query-mean led this early screen. Full family training later retained
decision-marker / option-close as the stable cross-backbone recipe.

\subsection{LoRA rank}

The two studies below use different bases, budgets, and selection rules.

\textbf{Qwen 4B early screen.} At the fixed causal checkpoint, aggregate ranking
over development and diagnostic Transfer metrics selected rank 32.

\begin{tabular}{p{.20\linewidth}rr|rr}
\toprule
 & \multicolumn{2}{c|}{Development} & \multicolumn{2}{c}{Transfer diagnostic} \\
Rank & Accuracy & NLL & Accuracy & NLL \\
\midrule
2  & 76.11 & 0.6385 & 67.59 & 0.9183 \\
4  & 76.33 & 0.6187 & 68.55 & 0.9013 \\
8  & 77.88 & 0.5405 & 70.08 & 0.8642 \\
16 & \good{78.87} & 0.5128 & 69.69 & 0.8601 \\
32 & 77.99 & \good{0.5010} & \good{70.27} & \good{0.7800} \\
\bottomrule
\end{tabular}

\textbf{27B base-start gate.} A separate 250-step, single-seed study compared
ranks 8, 16, 32, and 64. Rank 16 was the registered control and the only eligible
arm after multiplicity-corrected gates; rank 8 had better point estimates but did
not pass the adjusted promotion rule.

\begin{tabular}{p{.20\linewidth}rrrr}
\toprule
Rank & Dev. accuracy $\uparrow$ & Dev. NLL $\downarrow$ & Dev. Brier $\downarrow$ & Gate \\
\midrule
8  & \good{83.63} & \good{0.4445} & \good{0.2467} & not eligible \\
16 & 82.08 & 0.4708 & 0.2568 & \good{selected control} \\
32 & 82.08 & 0.4788 & 0.2599 & not eligible \\
64 & 82.52 & 0.4805 & 0.2579 & not eligible \\
\bottomrule
\end{tabular}

The production family ultimately uses rank 8 based on its separate full-run
cost/quality trade-off.

\subsection{Loss and calibration objectives}

The objective studies covered CE, pure InfoNCE, mixed CE--InfoNCE, supervised
Brier penalties, and RLCR-style policy terms.

\begin{itemize}
  \item In the controlled 27B base-start comparison, CE and pure InfoNCE had the
  same corrected development NLL to four decimals (0.4766), but InfoNCE worsened
  Brier from 0.2592 to 0.2700. CE was selected.
  \item In the early 4B sweep, contrastive weights from 0.01 through 3.0 were
  tested. A weight of 0.03 ranked first on the development aggregate, but its
  diagnostic Transfer accuracy was 67.97\%, below CE's 69.41\%. The production
  release therefore uses CE.
  \item A rank-16 factorial screen selected residual + CE + 0.1 contrastive weight
  within that local screen. Effects were small: under the residual head the
  contrastive term changed accuracy by $-0.11$ points and Brier by $-0.00157$.
  This did not override the later full-run recipe.
  \item Supervised Brier weights 0.1, 0.25, and 0.5 produced no arm satisfying the
  registered joint accuracy/calibration gate.
  \item Correctness-only, selected-action Brier, full-distribution Brier, and a
  calibration-penalized RLCR variant were separated in staged controls. No RLCR
  arm passed the complete promotion gate. In the earlier released SFT/RLCR pair,
  RLCR changed Transfer accuracy by $-0.10$ points (95\% paired interval
  $[-0.58,+0.39]$) while the small development-NLL gain did not survive
  independent recalibration.
\end{itemize}

\subsection{Data-mixture interventions}

The data studies changed broad task balance while holding the model and training
budget fixed. They are reported only at the aggregate level here.

\begin{itemize}
  \item Two targeted-text variants were tested against the same control at the
  fixed checkpoint. Variant v5 failed both calibrated-Brier improvement and
  accuracy non-inferiority; variant v6a improved its accuracy point estimate by
  about one point but still failed the calibrated-Brier promotion gate.
  \item A matched whole-mixture intervention and a matched hard-reasoning
  intervention were compared with their frozen control. Neither passed all
  registered gates, so no candidate was selected.
  \item A proposed dose study was stopped before training because no eligible
  fresh selection panel was available. Its status is a stopped protocol.
\end{itemize}

\subsection{Alternative decision architectures}

An exploratory route compared a jointly encoded cross scorer with dual-tower
alternatives on the same 1,046 Transfer decisions.

\begin{tabular}{p{.32\linewidth}rrrr}
\toprule
Architecture & Accuracy $\uparrow$ & NLL $\downarrow$ & Brier $\downarrow$ & ECE $\downarrow$ \\
\midrule
Cross-joint & \good{57.17} & \good{1.1766} & \good{0.5915} & 0.1494 \\
Scalar dual tower & 35.28 & 1.2459 & 0.6633 & \good{0.0136} \\
Expressive dual, scratch & 47.42 & 1.1958 & 0.6176 & 0.0717 \\
Expressive dual, official init & 41.78 & 1.2351 & 0.6390 & 0.0944 \\
\bottomrule
\end{tabular}

The cross-joint package beat the original dual tower, and the expressive scratch
dual improved substantially over the scalar dual. Nevertheless, the complete
registered gates failed: one dual comparator did not establish learning, and the
rescued variants did not jointly pass learning and cross-comparison conditions.
No architecture was promoted. These single-seed comparisons used packages
that differ from the released LoRA system.

A separate LM-token package compared the official next-token objective,
InfoNCE-soft, and hard/soft SoftCE variants. SoftCE-soft was promoted within that
local screen (43.69\% accuracy, 1.2627 NLL, 0.6501 Brier after calibration), but
the preregistration labels its comparison with JevAny as a same-benchmark system
comparison.

\subsection{Readout family and backbone scale}

The final release provides the cleanest practical readout comparison because the
two Qwen 4B models share the same base revision, data scale, rank, optimizer
family, and precision. Direct-token gains 0.86 JevBench points and improves all
three Transfer calibration metrics, while pointer gains 0.48 Transfer accuracy
points, supports more choices, and trains faster over a full trajectory.

Across released pointer models, scaling/backbone choice matters substantially:
Gemma 4B reaches 70.84\% Transfer, Qwen 4B 78.68\%, Muse 30B 83.46\%, and Qwen
27B 85.76\%. This family comparison varies architecture and parameter count
together.

\subsection{Ablation ledger}

\begin{tabular}{>{\raggedright\arraybackslash}p{.20\linewidth}>{\raggedright\arraybackslash}p{.27\linewidth}>{\raggedright\arraybackslash}p{.43\linewidth}}
\toprule
Track & Arms & Outcome \\
\midrule
Pointer head & Linear, MLP, residual & Residual selected in the controlled screen. \\
Representation & Marker/close, query mean, option mean, combined & Query mean led the early screen; marker/close retained for the full family. \\
LoRA rank & 2/4/8/16/32 at 4B; 8/16/32/64 at 27B & Local winners depended on protocol; production uses rank 8. \\
Contrastive loss & CE, InfoNCE, seven mixed weights & Small mixed terms helped some calibration screens; CE retained at full scale. \\
Supervised Brier & Weights 0.1/0.25/0.5 vs CE & No weight passed the joint gate. \\
Post-training & CE-path controls, correctness-only, selected/full Brier RLCR & No policy arm promoted; RLCR remains experimental. \\
Data mixture & Control, targeted variants, matched package/hardness changes & No candidate passed all gates; one planned dose study stopped pre-run. \\
Architecture & Cross-joint, scalar dual, expressive dual, initialized dual & Informative package differences, but complete promotion gates failed. \\
LM-token objective & Official CLM, InfoNCE-soft, SoftCE hard/soft & SoftCE-soft won its local screen; treated as a system-package result. \\
Decision readout & Pointer vs direct-token & Complementary 4B trade-off; both released. \\
Backbone family & Gemma, Qwen, Muse & Qwen3.8 27B is the highest-accuracy release. \\
\bottomrule
\end{tabular}

\section{Compute accounting}

Compute below covers training through each released checkpoint. GPU-hours equal
elapsed wall-clock hours times the actual DDP world size and exclude ablation,
evaluation, failed jobs, and training after the selected checkpoint.

\begin{tabular}{p{.34\linewidth}rrrr}
\toprule
Model & Step & Parallel H200s & Wall hours & H200 GPU-hours \\
\midrule
Gemma 4B LoRA & 2,771 & 32 & $\sim$3.28 & $\sim$104.9 \\
Qwen3.5 4B LoRA & 13,850 & 32 & 10.33 & 330.6 \\
Qwen3.5 4B Direct-Token & 9,695 & 32 & 8.41 & 269.1 \\
Muse Glimmer 30B LoRA & 3,324 & 40 & 2.88 & 115.4 \\
Qwen3.8 27B LoRA & 22,160 & 32 & 18.83 & 602.7 \\
\midrule
\textbf{Released-checkpoint total} & & & & \textbf{$\sim$1,423} \\
\bottomrule
\end{tabular}

Gemma is approximate because that earlier checkpoint format did not persist
cumulative elapsed seconds; its value uses the recorded run start and checkpoint
timestamp. Direct-token, Qwen 27B, and Muse use checkpoint-native elapsed time;
Qwen 4B pointer uses terminal trainer telemetry because the released checkpoint
is the completed step 13,850 run. Peak within-run parallelism was 40 GPUs.

\clearpage
\section{Reproducibility and evaluation scope}

\begin{summarybox}
\textbf{Release checks.} Base revisions are pinned. Every public artifact carries
its adapter metadata and tokenizer state. Release manifests, checkpoint-native
reload reports, and GPU loader/reconstruction parity were checked before
publication. Transfer covers 1,046/1,046 scored decisions and JevBench covers
231/231 public development items.
\end{summarybox}

\begin{itemize}
  \item Short screens use a single training seed unless stated otherwise.
  Measuring training-seed uncertainty requires repeated runs.
  \item Transfer served as a development diagnostic in several research stages
  and is reported on that basis in the release comparison.
  \item Backbone families differ in weights and may differ in architecture,
  tokenizer, and native media machinery.
  \item GPU-hours cover training through each selected checkpoint. Total research
  spend and inference cost are outside this accounting.
  \item The current release contains LoRA SFT checkpoints. Full-parameter SFT,
  improved post-training, and hyperparameter-optimized variants are roadmap
  items.
\end{itemize}

\section{Release map}

\begin{tabular}{p{.25\linewidth}p{.66\linewidth}}
\toprule
Artifact & Location \\
\midrule
Source code & \url{https://github.com/SimpleJev/JevAny} \\
Model collection & \href{https://huggingface.co/collections/tianxinwei/jevany-adaptive-decision-systems-6ab2c941bcecb4d2c61d1326}{Hugging Face / JevAny Adaptive Decision Systems} \\
Machine-readable metrics & \texttt{results/model-family-v2.json} \\
Algorithm reference & \texttt{docs/ALGORITHM.md} \\
Training guide & \texttt{docs/TRAINING.md} \\
Evaluation guide & \texttt{docs/EVALUATION.md} \\
\bottomrule
\end{tabular}

\vfill
\begin{center}
\color{muted}\small
Infrastructure adapted from Kev is
credited in the repository NOTICE and acknowledgements. Code is Apache-2.0;
base models and upstream datasets retain their own terms.
\end{center}

\end{document}
